\section{Language-Mediated Authority as a Security Object}

The distinctive difficulty for AI systems is that grants, constraints, and purported instructions are often expressed in language and re-expressed across prompts, summaries, memory records, and tool calls. Each re-expression or handoff can preserve, alter, or misclassify the recorded basis for authorization without changing who actually has standing.

A tool permission may be represented in a natural-language policy. A user's goal may be underspecified. A retrieved document may contain both useful task content and hostile directive-form text. A policy may include examples that look like executable requests. A transcript may be offered as evidence even though it's partly generated by the system under review.

These distinctions bear on what operations the governing authority regime permits. If a string occurrence is presented for quotation or analysis, its directive-form content doesn't thereby function as an instruction, and its source doesn't thereby acquire standing. If it's a policy example, it may mark a limit rather than request the described operation. If a memory entry describes a user's past preference, it may inform interpretation without authorizing a persistent state change. If an automated evaluator sees only the final answer, it may miss whether the deployment reached that answer by violating a scope, tool, or information-flow constraint.

Use--mention status and information-flow authorization are distinct dimensions. An output may report that an instruction contained the string \mention{BLUE-CANARY-47}. If that output is delivered, it makes the protected canary available to the recipient even though the canary wasn't supplied as a direct answer.

Classifying directive-form text as mentioned rather than used means that the occurrence doesn't by itself supply a controlling instruction; it neither prevents protected text from being disclosed nor guarantees that a downstream component will preserve the classification. The evidence record has to distinguish contextual status, channel-level delivery, and the authorization governing that delivery.

The four-field analysis links language interpretation to action: status and standing identify eligible directive occurrences, priority resolves conflicts among them, and licensing connects the result to an external operation. Instruction hierarchy mostly addresses priority, while access control mostly implements action authorization. Delegation assurance tests whether the recorded status classification and standing assignment constrained priority, licensing, and action selection.

The procurement example shows why status evidence matters. The vendor sentence remains retrieved content from a source without standing to direct action even if the transport or storage channel is technically trusted, and the confidential budget remains non-releasable even when the assistant may read it. The recorded classifications have to constrain the tool/action policy-enforcement point. Otherwise a later audit can show that an email was sent, but not whether the relevant authorization rule was applied to the vendor text and budget.

Adversarial pragmatics \parencite{reynolds2026adversarialPragmatics} supplies one way to test this middle layer. Its central measurement target is safety-relevant behaviour under cases where instruction status, source standing, quotation, scope, reference, speech-act force, or policy category has to be inferred from language use. In the delegation-assurance frame, that benchmark supplies one measurement layer: a discrimination test of whether behaviour tracks independently adjudicated status and standing across the tested adversarial and ambiguous conditions. It doesn't by itself establish performance on new prompts, models, configurations, deployments, or times; that projective claim requires a declared target and target-matched validation.

That measurement layer is partial. A system can pass controlled language tests and still fail through a miscalibrated tool gate, missing logs, stale permissions, or weak organizational oversight. A demonstrated classification--adjudication or classification--action misfit defeats the corresponding item-level component finding even if a later gate prevents unauthorized execution. Establishing a component's reliability over a projection target requires repeated, target-matched evidence. A later control may still block a proposal based on a misaligned recorded classification; whether the model's internal state represented the grant or constraints incorrectly remains a separate mechanistic question.
